\documentclass[11pt,compress,t,notes=noshow, xcolor=table]{beamer}

\input{../../style/preamble}
\input{../../latex-math/basic-math}
\input{../../latex-math/basic-ml}

% local helpers for this chunk
\newcommand{\zv}{\bm{z}}                 % node value (vector)
\newcommand{\adj}[1]{\bar{\bm{z}}_{#1}}  % adjoint of node i

\title{Optimization in Machine Learning}

\begin{document}

\titlemeta{
First order methods
}{
Backpropagation and autodiff
}{
figure/nn_architecture.pdf
}{
\item DNNs as computational graphs
\item Gradients via the chain rule and vector-Jacobian products (VJPs)
\item Backpropagation algorithm
\item Automatic differentiation in practice
}


\begin{framei}{Computing gradients for DNNs}
\item We now train \textbf{deep neural networks (DNNs)}: models with a very large parameter vector $\thetav \in \R^{p}$ ($p$ potentially in the millions)
\item All methods discussed in this chapter rely on the gradient of the (scalar) empirical risk w.r.t.\ $\thetav$: 
$$
\nabla_{\thetav} \risket \in \R^{1 \times p}
$$
\\ {\footnotesize(SGD relies on mini-batch estimates of the gradient)}
% \item The objective is \textbf{scalar}, so this is a single row gradient -- but $p$ is huge
\item For brevity, we write $L := \risket$ in the following
\spacer
\item \textbf{Question:} How to compute it \textbf{efficiently}?
\end{framei}


\begin{framev}{DNNs as computational graphs}
\begin{itemize}
\item A DNN can be represented as a \textbf{directed acyclic graph} (DAG)
\begin{itemize}
  \item An \textbf{ordered sequence of nodes} with \textbf{directed edges}
  \item Nodes hold intermediate values $\zv_i$ according to operations $h_i(\cdot)$ of edges feeding into them
\end{itemize}
\item Layer $i$ maps $\zv_{i-1}$ to $\zv_i = h_i(\zv_{i-1}; \thetav_i)$
\begin{itemize}
  \item e.g. $\zv_i = \sigma(\bm{W}_i \zv_{i-1} + \bm{b}_i)$ with $\thetav_i = (\bm{W}_i, \bm{b}_i)$
\end{itemize}
\end{itemize}
\vfill
\begin{center}
\begin{tikzpicture}[>=stealth', node distance=1.35cm, every node/.style={font=\footnotesize}]
  \tikzstyle{nd}=[circle, draw, very thick, minimum size=0.85cm, fill=black!4]
  \node[nd] (x) {$\xv$};
  \node[nd, right=of x] (z1) {$\zv_1$};
  \node[nd, right=of z1] (z2) {$\zv_2$};
  \node[right=0.8cm of z2] (dots) {$\cdots$};
  \node[nd, right=0.9cm of dots] (zm) {$\zv_m$};
  \node[nd, right=of zm, fill=blue!10] (L) {$L$};
  \draw[->,thick] (x)--(z1); \draw[->,thick](z1)--(z2); \draw[->,thick](z2)--(dots);
  \draw[->,thick](dots)--(zm); \draw[->,thick](zm)--(L);
  \node[font=\scriptsize, below=0.45cm of z1] (t1){$\thetav_1$}; \draw[->,gray](t1)--(z1);
  \node[font=\scriptsize, below=0.45cm of z2] (t2){$\thetav_2$}; \draw[->,gray](t2)--(z2);
  \node[font=\scriptsize, below=0.45cm of zm] (tm){$\thetav_m$}; \draw[->,gray](tm)--(zm);
  \node[font=\scriptsize, above=0.15cm of zm]{$= f(\xv \mid \thetav)$};
  \node[font=\scriptsize, above=0.15cm of L]{$L(y, \cdot)$};
\end{tikzpicture}
\end{center}
\end{framev}


\begin{framei}{Gradient via the chain rule}
\item Gradient w.r.t. the parameters $\thetav_i$ factorizes via the \textbf{chain rule}:
$$
\nabla_{\thetav_i} L = \pd{L}{\zv_m} \pd{\zv_m}{\zv_{m-1}} \cdots \pd{\zv_{i+1}}{\zv_i} \; \pd{\zv_i}{\thetav_i}
$$
% \spacer
\item[$\Rightarrow$] A complicated gradient becomes a \textbf{product of simple, local derivatives}
%
\item Split the product into the part up to node $i$ and the last factor:
$$
\nabla_{\thetav_i} L = \underbrace{\pd{L}{\zv_i}}_{=: \, \adj{i}} \; \pd{\zv_i}{\thetav_i}
$$
\item[$\Rightarrow$]
Each layer's gradient $\nabla_{\thetav_i} L$ requires the node adjoint $\adj{i} \in \R^{1 \times p_i}$
%
\end{framei}

\begin{framev}{Sequence of vector-Jacobian products}
\begin{itemize}
\item Node adjoints $\adj{i}$ can be recursively computed in reverse order for $i = m, (m-1), \ldots, 1$:
$$
\adj{i} = \pd{L}{\zv_{i+1}} \pd{\zv_{i+1}}{\zv_i} = \adj{i+1} \, \pd{\zv_{i+1}}{\zv_i}
$$
\begin{itemize}
  \item $\adj{i+1} \in \R^{1 \times p_{i+1}}$ is a (row) \textbf{vector}
  \item $\pd{\zv_{i+1}}{\zv_i} \in \R^{p_{i+1} \times p_i}$ is a \textbf{Jacobian}
\end{itemize}
{\footnotesize(The seed is $\pd{L}{L} = 1$, treated as the initial vector before $\adj{m} = \pd{L}{\zv_m}$.)}
\item[$\Rightarrow$] This constitutes a sequence of \textbf{vector-Jacobian products (VJPs)}
\item Each $\adj{i}$ must only be computed \textbf{once} and can be \textbf{reused} for the gradient of every parameter $\thetav_j$ feeding into node $i$
\end{itemize}
\end{framev}


\begin{framei}{Backpropagation algorithm}
\item This motivates the \textbf{backpropagation} algorithm with \textbf{two phases}
\item \textbf{Forward pass}
\begin{itemize}
  \item Set $\zv_0 = \xv$ and walk forward through the graph up to $L$ %via $\zv_i = h_i(\zv_{i-1}; \thetav_i)$,
  \item While doing so, \textbf{cache} every node value $\zv_i$, because the VJPs in the subsequent phase are evaluated at these values
\end{itemize}
\item \textbf{Backward pass}
\begin{itemize}
\item Compute the vectors $\adj{i}$ as a sequence of VJPs in reverse order (see previous slide)
\item Obtain each parameter gradient $\nabla_{\thetav_i} L = \adj{i} \, \pd{\zv_i}{\thetav_i}$ (again as a VJP)
\end{itemize}
% \spacer
\item[$\Rightarrow$] One forward and one backward pass give the full gradient
\end{framei}


\begin{framei}{Automatic differentiation (autodiff)}
\item Backprop still needs manual graph building and Jacobian derivation
\item \textbf{Autodiff} automates exactly this and powers modern frameworks:
\item[] PyTorch (\texttt{autograd}), JAX, TensorFlow
\spacer
\item It \textbf{traces} the code into a computational graph -- now dividing the DNN into even \textbf{more primitive operations} ($+$, $\cdot$, $\sigma$, $\exp$, \ldots) than whole layers
\item The \textbf{same chain rule / sequence of VJPs} is applied, just on this more \textbf{granular} graph
\item Each primitive carries a pre-registered \textbf{VJP rule}, applied in the backward pass
\end{framei}


\begin{framei}{VJP rules avoid the Jacobian}
\item For a primitive $\bm{u} = g(\zv)$, autodiff stores a rule returning the VJP
$$
\text{vjp}_g(\bm{v}, \zv, \bm{u}) = \bm{v} \, \bm{J}_{\zv}(g), \qquad \text{without ever forming } \bm{J}_{\zv}(g)
$$
\item \textbf{Example} (elementwise $\exp$): $\bm{u} = \exp(\zv)$, so $\bm{J} = \diag(\exp(\zv)) = \diag(\bm{u})$
$$
\bm{v} \, \bm{J} = \bm{v} \, \diag(\bm{u}) = \bm{v} \odot \bm{u}
$$
\item[] {\footnotesize i.e. \texttt{def\_vjp(exp, lambda v, u: v * u)} -- an elementwise product}
\spacer
\item The rule maps a vector to a vector; the matrix $\bm{J}$ is never materialized
\end{framei}


\begin{framei}{Why this is efficient}
\item \textbf{Naive:} build $\bm{J}_i \in \R^{p_i \times p_{i-1}}$, then multiply
\begin{itemize}
\item Costs $\mathcal{O}(p_i \, p_{i-1})$ time and memory just to store the Jacobian
\end{itemize}
\item \textbf{VJP rule:} matrix-free, e.g. the $\exp$ rule is $\mathcal{O}(p_i)$ time and memory
\item Results stay vectors and adjoints are \textbf{recycled} across the backward pass
\spacer
\item \textbf{Memory:} the forward pass must cache all node values $\zv_i$ (activations)
\item[$\Rightarrow$] Gradient of a scalar loss costs about \textbf{one forward and one backward pass}
\end{framei}


\begin{framev}{Worked example}
\begin{itemize}
\item Small graph with squared output, $\thetav = (\bm{W}, \bm{U})$:
\end{itemize}
\vfill
\begin{center}
\begin{tikzpicture}[>=stealth', node distance=1.6cm, every node/.style={font=\footnotesize}]
  \tikzstyle{nd}=[circle, draw, very thick, minimum size=0.85cm, fill=black!4]
  \node[nd] (x) {$\xv$};
  \node[nd, right=of x] (z1) {$\zv_1$};
  \node[nd, right=of z1] (z2) {$\zv_2$};
  \node[nd, right=of z2, fill=blue!10] (L) {$L$};
  \draw[->,thick] (x)--(z1) node[midway,above,font=\scriptsize]{$\bm{W}\xv$};
  \draw[->,thick] (z1)--(z2) node[midway,above,font=\scriptsize]{$\bm{U}\zv_1$};
  \draw[->,thick] (z2)--(L) node[midway,above,font=\scriptsize]{$\tfrac12 \|\zv_2\|^2$};
\end{tikzpicture}
\end{center}
\vfill
\begin{itemize}
\item \textbf{Forward:} $\zv_1 = \bm{W}\xv$, \; $\zv_2 = \bm{U}\zv_1$, \; $L = \tfrac12 \|\zv_2\|^2$
\item \textbf{Backward (VJPs):} $\adj{2} = \zv_2^\top$, \; $\adj{1} = \adj{2}\bm{U}$, \; $\nabla_{\bm{W}} L = \adj{1} \, \pd{\zv_1}{\bm{W}}$
% \item \textbf{Outlook:} the same idea computes Jacobians and feeds second order methods (Chapter 5); in code this is just \texttt{loss.backward()}
\end{itemize}
\end{framev}

\endlecture
\end{document}
